\subsection{Defnition}
The idea of the trajectory controller is to minimize the error between target height and real max height.
To do so we try to minimize following error throghout the flight
\begin{equation}
    e = h_f - h,
\end{equation}
where $h_f$ is target height and $h$ is the predicted final (apogee) altitude from the current state --
NOT the altitude reached so far.

The idea is that controller input is the final height error and output the air brake extension percentage,
writing it as
\begin{equation}
    \frac{\text{ext}(s)}{e(s)} = C_1(s).
\end{equation}

The loop then consists of air brake extension controller $C_2(s)$, plant $P(s)$\footnote{The $s$-domain
notation used throughout this section is block-diagram/signal-flow shorthand, not literal LTI transfer
functions: the predictor ($R(s)$ below) and the plant ($P(s)$) are both nonlinear (the predictor through
$\log$/$v^2$, the plant through quadratic drag and, later, attitude dynamics), and the real system is
sampled (fixed \texttt{sample\_rate\_hz}, zero-order hold), so $z$-domain would be the more honest notation
once past this purely structural diagram. A constant gain like $C_1(s)=k_p$ happens to be trivially valid
as written; $R(s)$ and $P(s)$ are not.}, which gives the current
state of the rocket, which then is fed into height prediction alogirthm

\subsubsection{P/PD - Burnout - 1}
At the begining we will assume that air brake drag is linear, of form
\begin{equation}
    C_d = C_{d,f}\cdot f,
\end{equation}
there is no intertia and controller of the air brake extension, i.e. the movement is instatanous, and
the point mass plant is used.

The system's logic is that using the results from plant, the predictor can estimate final altitude, which
can then be fed into error calculation $e = h_f - h$, which then is being used by controller, which
controls the extension $f$, which is fed into the plant and the loop closes. The system can be seen on
\ref{fig:1}

Starting with predictor at the beggining the analytical derivation of after burnout dynamics is used,
it follows that 
\begin{equation}
    m \frac{dv}{dt} = -mg - kv^2.
\end{equation}
By applying approximation
\begin{equation}
    \frac{dv}{dt} = \frac{dv}{dt}\frac{dh}{dh} = \frac{dh}{dt}\frac{dv}{dh} = v \frac{dv}{dh}
\end{equation}
and $u = v^2$. One gets
\begin{equation}
    \frac{du}{dh} = -2g -\frac{2k}{m}u.
\end{equation}
This intern can be integrated
\begin{equation}
    \int_{u}^{u_f} \frac{du}{-2g -\frac{2k}{m}u} = \int_{h}^{h_f}dh
\end{equation}
where $u = v^2$ at $h$ and $u_f = 0$ and $h_f$, gives sollution of error
\begin{equation}
    e = h_f - h = \frac{m}{2k}\ln\left(1 + \frac{k v^2}{mg}\right),
\end{equation}
which as output gives the error. 

This predictor transfer function is of form
\begin{equation}
    R(s) = \frac{e(s)}{x^*(s)},
\end{equation}
where it takes vector of chosen states (for now $h, v, f$) and outputs the error. 
It's worth noting that for this case the boundary conditions ($h_t, v_t$) are a side input.
\todo{write about how such system can be translatet into velocity based trajectory; idk the reason
but for example not to reach a transonic region or smn}

The error is then fed up to the P controller, which by definition is given as
\begin{equation}
    C_1(s) = \frac{f(s)}{e(s)} = k_p.
\end{equation}

The brake extension is then, given to a plant
\begin{equation}
    P(s) = \frac{x(s)}{f(s)},
\end{equation}
which returns a vector of states. This vector of state is then fed back in the predictor and the loop closes.

\todo{results}





